\subsection*{最初に必要な用語}
\addcontentsline{toc}{subsection}{最初に必要な用語}
構成管理では略語が多い。初学者は、まず「構成品目」「基準版」「変更要求」「状態記録」「監査」の五つを軸に読むとよい。

\par\smallskip\noindent\refstepcounter{table}\label{tab:ch08-01}表 \thetable: 最初に必要な用語における用語・初学者向けの説明\par\smallskip
表 \thetable は、最初に必要な用語における用語・初学者向けの説明を比較・確認しやすい形で整理したものである。\par\smallskip
\begin{longtable}{p{0.30\linewidth}p{0.64\linewidth}}\toprule
用語 & 初学者向けの説明 \\ \midrule
\endfirsthead\toprule
用語 & 初学者向けの説明 \\ \midrule
\endhead
\term{CR}{Change Request} & 変更要求。プロジェクト範囲、費用、計画、コード、日程などの変更を求める依頼である。 \\
\term{CCB / SCCB}{Configuration Control Board / Software Configuration Control Board} & 構成管理委員会 / ソフトウェア構成管理委員会。変更を承認、修正、延期、却下する権限を持つ。 \\
CI & 構成品目。ひとつの管理単位として扱う品目である。 \\
\term{SCI}{Software Configuration Item} & ソフトウェア構成品目。ソフトウェアに関する構成品目である。 \\
\term{CM / SCM}{Configuration Management / Software Configuration Management} & 構成管理 / ソフトウェア構成管理。品目を識別し、変更を統制し、状態を記録し、監査する管理活動である。 \\
\term{SCMP}{Software Configuration Management Plan} & ソフトウェア構成管理計画。構成管理をどう進めるかを定める生きた文書である。 \\
\term{SCSA}{Software Configuration Status Accounting} & ソフトウェア構成状態記録・報告。構成品目、基準版、変更、関係を記録し報告する活動である。 \\
\term{FCA / PCA}{Functional Configuration Audit / Physical Configuration Audit} & 機能構成監査 / 物理構成監査。仕様に合うか、実物と文書が合うかを確認する。 \\
\term{SBOM}{Software Bill of Materials} & ソフトウェア部品表。ビルドに使われた部品と供給網上の関係を記録する。 \\
\term{VDD}{Version Description Document} & 版記述文書。リリースに何が含まれるかを物理的に記述する文書である。 \\
\bottomrule\end{longtable}
